\documentclass[12pt]{article}
\usepackage[utf8]{inputenc}
\usepackage[T1]{fontenc}
\usepackage{hyperref}
\title{Wia Language Resources for K-12 Curriculum Development}
\author{Compiled by Claude Sonnet 5.5 (claude-sonnet-5-5) for Andr\'as Kornai}
\date{}
\begin{document}
\maketitle

\section{Introduction}
Wia (ISO 639-3 ssj) is an endangered language of Madang Province, Papua New Guinea, with about 93 speakers. Wikipedia lists it as Sausi, a Rai Coast language, also known as Uya. Joshua Project gives alternate names including Witi and Wiya. A separate Papuan language called Tabo or Waia is spoken in Western Province. Do not confuse it with this one. The web search found very little for this language, and the connection between the names Wia and Sausi rests on a single source (Joshua Project). It should be verified against Ethnologue or Glottolog.

\section{Primary Educational Resources}
No primers, textbooks or school materials in Wia turned up in the search.

\section{Linguistic and Reference Materials}
\begin{itemize}
\item Wikipedia, Sausi language (overview, speaker count): \url{https://en.wikipedia.org/wiki/Sausi_language}
\end{itemize}

\section{Religious and Community Resources}
\begin{itemize}
\item Joshua Project, language page for ssj. It lists alternate names and reports Scripture translation as started: \url{https://joshuaproject.net/languages/ssj}
\item Joshua Project, Sausi, Uya people group. It describes the community as under 1,000 people near Koropa, upper Ramu River, Madang: \url{https://joshuaproject.net/people_groups/14789}
\end{itemize}

\section{Digital and Audio Resources}
No audio, video or software resources turned up in the search.

\section{Curriculum Development Context}
Almost nothing is available in the search results. A curriculum would need new documentation work: a wordlist, an orthography and recordings, made together with the community. Outside Wikipedia and Joshua Project, any fieldwork or archive holdings would have to be found separately.
\end{document}